\documentclass[12pt]{article}
\usepackage[utf8]{inputenc}
\usepackage[T1]{fontenc}
\usepackage[margin=1in]{geometry}
\usepackage{setspace}
\usepackage{xcolor}
\usepackage[colorlinks=true, urlcolor=blue!60!black, linkcolor=blue!60!black]{hyperref}
\setlength{\parindent}{2em}
\setlength{\parskip}{1.25em}
\renewcommand{\baselinestretch}{1.0}
\begin{document}
\begin{center}
{\LARGE \textit{West Virginia v. EPA}, 597 U.S. 697 (2022)\par}
\vspace{0.8em}
{\large Case Brief\par}
\end{center}
\vspace{1em}

\section{Memory Jogger}
Giving the ``major questions doctrine'' its name in a majority opinion for the first time, the Court held 6--3 that \S~111(d) of the Clean Air Act did not clearly authorize EPA's Clean Power Plan, which set carbon-dioxide limits premised on shifting electricity generation away from coal (\href{https://supreme.justia.com/cases/federal/us/597/20-1530/}{Justia}).

\section{Detailed Case Facts}
The Clean Air Act (CAA) establishes three primary stationary-source programs: the National Ambient Air Quality Standards (NAAQS) program (\S\S~108--110, 42 U.S.C. \S\S~7408--7410), the Hazardous Air Pollutants (HAP) program (\S~112, \S~7412), and the New Source Performance Standards program (\S~111, \S~7411). Section 111(d) operates as a ``gap-filler'' requiring the Environmental Protection Agency (EPA) to set emissions guidelines for \textit{existing} stationary sources emitting pollutants regulated for new sources under \S~111(b) but not covered by the NAAQS or HAP programs. Under \S~7411(a)(1), EPA determines the ``best system of emission reduction'' (BSER) that is ``adequately demonstrated'' (taking into account cost, nonair quality health and environmental impacts, and energy requirements) and quantifies the achievable emissions limit; States then submit plans enforcing that limit. Prior to 2015, EPA had used \S~111(d) only a handful of times since 1970, always basing BSER on facility-level, technology-based controls (such as scrubbers or fiber mist eliminators) that cause regulated plants to operate more cleanly.

In October 2015, the Obama-era EPA promulgated the Clean Power Plan (CPP) under \S~111(d) to curb carbon dioxide emissions from existing coal- and natural-gas-fired power plants. EPA's BSER consisted of three ``building blocks'': (1) ``heat rate improvements'' at coal-fired plants, (2) ``generation shifting'' from coal-fired to natural-gas-fired plants, and (3) ``generation shifting'' from coal and gas plants to renewable sources (wind and solar). Operators could comply by reducing their own coal generation, investing in cleaner plants, or purchasing credits in a cap-and-trade regime. EPA projected that coal's share of national electricity generation would drop from 38\% in 2014 to 27\% by 2030, and set emissions ceilings so strict that no existing coal plant could comply through on-site technology alone.

\section{Procedural History}
The same day EPA promulgated the CPP, 27 States and dozens of other parties petitioned for review in the D.C. Circuit. After the D.C. Circuit denied a stay, the Supreme Court granted an unprecedented stay preventing the CPP from taking effect (\textit{West Virginia v. EPA}, 577 U.S. 1126 (2016)). Following the 2016 presidential election, the Trump-era EPA repealed the CPP in 2019---concluding that \S~111(d) limited BSER to source-specific, inside-the-fenceline measures and that the major questions doctrine precluded generation shifting---and replaced it with the narrower Affordable Clean Energy (ACE) Rule, which relied solely on heat-rate improvements.

States and environmental groups petitioned for review in the D.C. Circuit, and West Virginia, North Dakota, Westmoreland Mining Holdings LLC, and The North American Coal Corporation intervened to defend the repeal and ACE Rule. On January 19, 2021, the D.C. Circuit held that EPA's repeal rested on a mistaken reading of the CAA, rejected application of the major questions doctrine, and vacated both the CPP repeal and the ACE Rule (\textit{American Lung Assn. v. EPA}, 985 F.3d 914 (D.C. Cir. 2021)). Following another change in administrations, EPA moved to stay the D.C. Circuit's mandate as to the CPP repeal while it pursued a new \S~111(d) rulemaking, and the court agreed. West Virginia, North Dakota, and the coal-industry intervenors petitioned for certiorari, which the Supreme Court granted and consolidated.

\section{Judicial Votes}
The Supreme Court \textbf{reversed and remanded} by a \textbf{6--3} vote. The case was argued on February 28, 2022, and decided on June 30, 2022 (\href{https://supreme.justia.com/cases/federal/us/597/20-1530/}{Justia}; \href{https://www.law.cornell.edu/supremecourt/text/20-1530}{Cornell LII}):

\begin{itemize}
    \item \textbf{Majority Opinion:} Chief Justice Roberts delivered the opinion of the Court, joined by Justices Thomas, Alito, Gorsuch, Kavanaugh, and Barrett.
    \item \textbf{Concurring Opinion:} Justice Gorsuch filed a concurring opinion, joined by Justice Alito.
    \item \textbf{Dissenting Opinion:} Justice Kagan filed a dissenting opinion, joined by Justices Breyer and Sotomayor.
\end{itemize}

\section{Holding}
\textbf{Questions Presented:}

\begin{enumerate}
  \item \textbf{Justiciability:} Whether state petitioners retained Article III standing and a live, non-moot controversy to challenge the D.C. Circuit's vacatur of the CPP repeal when EPA represented it did not intend to enforce the CPP and planned a new rulemaking.
  \item \textbf{Merits:} Whether \S~111(d) of the Clean Air Act, 42 U.S.C. \S~7411(d), authorized EPA to devise carbon-dioxide emissions caps based on the grid-wide ``generation shifting'' approach adopted in the Clean Power Plan.
\end{enumerate}

\textbf{Holding:} Reversed and remanded (6--3). First, the case remained justiciable: the D.C. Circuit's judgment vacating the ACE Rule and its embedded CPP repeal purported to bring the CPP back into legal effect, injuring the petitioner States as the direct objects of its regulatory requirement, and EPA's voluntary cessation did not moot the dispute because the Government vigorously defended generation shifting and refused to disclaim reimposing it. Second, on the merits, Congress did not grant EPA authority in \S~111(d) to devise emissions caps based on the CPP's generation-shifting approach.

\textbf{Rule (Major Questions Doctrine):} In ``extraordinary cases'' where the ``history and the breadth of the authority'' asserted by an agency and the ``economic and political significance'' of that assertion provide a ``reason to hesitate before concluding that Congress'' meant to confer such power, a merely plausible textual basis is insufficient; the agency must point to ``clear congressional authorization'' for the authority it claims.

\section{Analysis of Opinions}
\subsection{Majority Opinion (Chief Justice Roberts)}

\textbf{1. Formalizing the Major Questions Doctrine.} Chief Justice Roberts synthesized prior precedents---\textit{FDA v. Brown \& Williamson Tobacco Corp.}, 529 U.S. 120 (2000), \textit{Gonzales v. Oregon}, 546 U.S. 243 (2006), \textit{Utility Air Regulatory Group v. EPA}, 573 U.S. 302 (2014), \textit{Alabama Assn. of Realtors v. HHS}, 594 U.S. 758 (2021), and \textit{NFIB v. OSHA}, 595 U.S. 109 (2022)---explicitly adopting the label ``major questions doctrine'' for the first time in a majority opinion. Grounded in both separation-of-powers principles and a practical understanding of legislative intent, the doctrine presumes that Congress intends to make major policy decisions itself rather than leaving them to agencies through modest words, vague terms, or subtle devices.

\textbf{2. Why the Clean Power Plan Triggered the Doctrine.} The Court identified four converging indicators:
\begin{itemize}
  \item \textbf{Unheralded Power in an Ancillary Provision:} EPA claimed to discover a ``transformative expansion'' of regulatory authority in \S~111(d), a rarely used ``gap-filler'' and statutory ``backwater'' that for four decades had been understood as a source-specific, technology-based program (distinguishing the 2005 Mercury Rule because its cap was set by achievable wet-scrubber technology).
  \item \textbf{Fundamental Statutory Revision:} Shifting from making individual power plants operate more cleanly to forcing a sector-wide transition away from coal effected a ``fundamental revision of the statute,'' potentially empowering EPA to force coal plants to cease operating altogether.
  \item \textbf{Lack of Comparative Expertise:} Balancing grid reliability, electricity transmission, distribution, storage, and energy prices lies outside EPA's traditional environmental expertise---a point EPA itself had acknowledged in appropriations requests.
  \item \textbf{Conspicuous Congressional Rejection:} The CPP effectively enacted a carbon cap-and-trade regime that Congress had ``considered and rejected'' numerous times (e.g., the 2009 American Clean Energy and Security Act), making an oblique delegation through \S~111(d) all the more suspect.
\end{itemize}

\textbf{3. Absence of Clear Congressional Authorization.} Turning to the statutory text, the Court held that the word ``system'' in ``best system of emission reduction,'' shorn of context, is an ``empty vessel'' that cannot supply clear congressional authorization. Although the Acid Rain and NAAQS programs reference cap-and-trade ``systems,'' those provisions authorize trading only as a compliance mechanism to meet an emissions cap already set by Congress or objective ambient health standards---not to empower EPA to choose the cap itself. Finally, the Court emphasized the narrowness of its holding: it did not decide whether \S~111(d) is exclusively limited to inside-the-fenceline technological controls, holding only that the CPP's generation-shifting BSER exceeded \S~111(d).

\subsection{Concurring Opinion (Justice Gorsuch)}

Justice Gorsuch framed the major questions doctrine explicitly as a substantive ``clear-statement rule'' enforcing Article I's Vesting Clause and the nondelegation principle---parallel to clear-statement rules protecting sovereign immunity, federalism, and non-retroactivity. Gorsuch identified three triggers for a major question (great political significance / failed congressional proposals; regulation of a significant portion of the economy or billions in spending; and intrusion into traditional state-law domains such as utility regulation) and four clues for whether Congress spoke clearly (place in the overall statutory scheme, age and focus of the statute, past agency interpretations, and mismatch with agency expertise).

\subsection{Dissenting Opinion (Justice Kagan)}

Justice Kagan criticized the majority for issuing an advisory opinion on a defunct rule whose emissions targets market forces had already surpassed, and for abandoning textualism by unveiling the ``major questions doctrine'' as a ``get-out-of-text-free card.'' Textually, ``best system of emission reduction'' is broad rather than vague; Congress deliberately omitted the word ``technological'' for existing sources in the 1977 and 1990 Amendments; and cap-and-trade is an ordinary pollution-control ``system'' squarely within EPA's wheelhouse (\textit{American Elec. Power Co. v. Connecticut}, 564 U.S. 410 (2011)). Kagan argued that prior cases like \textit{Brown \& Williamson} were ordinary text-in-context interpretation where an agency strayed outside its lane and clashed with statutory structure, whereas here Congress deliberately delegated broad, adaptive authority to an expert agency to confront climate change.

\subsection{Subsequent Developments}

\begin{itemize}

\item \textbf{\href{https://supreme.justia.com/cases/federal/us/600/22-506/}{\textit{Biden v. Nebraska}, 600 U.S. 477 (2023)}:} In a 6--3 opinion by Chief Justice Roberts, the Court applied \textit{West Virginia}'s major questions framework to invalidate the Secretary of Education's \$430 billion student-loan cancellation program under the HEROES Act's authority to ``waive or modify'' student-loan provisions during a national emergency. Crucially, Justice Barrett concurred separately to reject Justice Gorsuch's view of the doctrine as a strong-form substantive ``clear-statement'' canon imposing a ``clarity tax'' on Congress; instead, Barrett reframed it as an ordinary contextual-textualist tool reflecting ``common sense'' and constitutional structure (illustrating the point with a babysitter handed a credit card and told to ``make sure the kids have fun'').

\item \textbf{\href{https://supreme.justia.com/cases/federal/us/607/24-1287/}{\textit{Learning Resources, Inc. v. Trump}, 607 U.S. 229 (2026)}:} Decided February 20, 2026, \textit{Learning Resources} provided the defining bipartisan test of \textit{West Virginia}'s legacy. In a 6--3 ruling, the Court held that the International Emergency Economic Powers Act (IEEPA)---authorizing the President to ``regulate \dots importation'' during a declared national emergency---did \textit{not} authorize President Trump's sweeping drug-trafficking and ``reciprocal'' tariffs. The opinions exposed a four-way split over \textit{West Virginia}:
    \begin{enumerate}
      \item \textit{Roberts Plurality (joined in Part II-A-2 only by Gorsuch and Barrett):} Applied \textit{West Virginia}'s major questions doctrine to IEEPA, holding that the President lacked ``clear congressional authorization'' to exercise Congress's core Article I taxing/tariff power and rejecting both emergency-statute and foreign-affairs carve-outs (``There is no major questions exception to the major questions doctrine'').
      \item \textit{Gorsuch vs. Barrett Concurrences:} Justice Gorsuch defended the major questions doctrine as a ``pro-Congress,'' common-law-rooted substantive clear-statement rule protecting Article I, while Justice Barrett reiterated her \textit{Nebraska} view that the doctrine is ordinary contextual textualism and rejected Gorsuch's strong-form ``clarity tax.''
      \item \textit{Kagan Concurrence (joined by Sotomayor and Jackson):} Agreed that IEEPA did not authorize tariffs under ``straight-up statutory construction'' (reading ``regulate'' in context alongside IEEPA's other eight verbs and Title 19's explicit tariff statutes), explicitly declining to join Part II-A-2 because ``I need no major-questions thumb on the interpretive scales.''
      \item \textit{Kavanaugh Dissent (joined by Thomas and Alito) and Thomas Dissent:} Justice Kavanaugh accepted the major questions doctrine in domestic cases but argued that IEEPA supplied clear authority and that a foreign-affairs exception should govern tariffs. Justice Thomas separately argued that nondelegation applies only to ``core legislative power'' governing deprivations of life, liberty, or property, and not to foreign-commerce import duties.
    \end{enumerate}
\end{itemize}

\section{Examples: Future Applications}
\begin{enumerate}

\item \textbf{Same-Side (Rule Likely Invalid): Repurposing a 1974 Biosecurity Statute to License Frontier AI.} In 1974, Congress enacted the Federal Biosecurity Preparedness Act, authorizing the Secretary of Health and Human Services to ``promulgate such standards and systems as the Secretary determines necessary to prevent catastrophic biological threats.'' For fifty years, HHS used the statute only to set lab-containment protocols for federal research facilities. Following a major public debate in which Congress repeatedly failed to pass a bill regulating artificial intelligence, the Secretary invokes the 1974 Act to issue a rule requiring every U.S. technology company training a frontier AI model above a specified compute threshold to obtain a federal license and submit model weights to HHS for biosecurity red-teaming, at an estimated industry compliance cost of \$18 billion per year. Challengers have a textbook major-questions claim: (1) an unheralded, transformative expansion of authority discovered in a 50-year-old statute used only for physical labs; (2) vast economic (\$18 billion/year) and political significance; (3) a stark mismatch between HHS's public-health mission and comprehensive licensing of the frontier AI sector; and (4) conspicuous congressional failure to enact AI licensing legislation. Under Roberts/Gorsuch, ``standards and systems'' is an ``empty vessel'' lacking clear congressional authorization; under Barrett/Kagan, even ordinary text-in-context interpretation makes it implausible that Congress hid a national AI-licensing regime inside a 1974 lab-safety clause.

\item \textbf{Same-Side (Rule Likely Invalid): Using Ratemaking Authority to Phase Out Coal.} Relying on its decades-old authority to ensure that wholesale electricity rates are ``just and reasonable,'' a federal energy commission declares that any rate charged by a coal-fired plant after 2035 is unjust and unreasonable, effectively requiring every such plant to retire. As in \textit{West Virginia}, the agency would be using a long-standing, general provision to restructure the Nation's electricity mix---a decision of vast economic and political significance that Congress has debated without acting. Absent clear congressional authorization, a court applying the major questions doctrine would likely hold that the ratemaking provision does not reach so far (\href{https://supreme.justia.com/cases/federal/us/597/20-1530/}{Justia}).

\item \textbf{Opposite-Side (Rule Likely Valid Under the Statute): Clear Congressional Authorization for AI Licensing.} Return to the biosecurity hypothetical, but suppose Congress in 2025 enacted the Artificial Intelligence Safety Act expressly granting a newly created AI Safety Agency authority to ``license and audit frontier AI models exceeding $10^{26}$ FLOPs to mitigate chemical, biological, radiological, or cyber risks,'' and the Agency sets detailed red-teaming standards under that statute. Here, the \textit{major questions doctrine does not defeat the rule}: Congress spoke clearly, the action sits squarely in the agency's core mission and expertise, and the statute is contemporaneous rather than an old ``backwater.'' Any challenge must therefore shift from statutory major-questions review to a constitutional \textit{nondelegation} claim---which, after \textit{FCC v. Consumers' Research}, 606 U.S. 656 (2025), will fail so long as Congress supplied an ``intelligible principle'' setting the general policy and boundaries of delegated authority.

\item \textbf{Opposite-Side (Doctrine Likely Not Triggered): A Traditional, Source-Specific Standard Under \S~111(d).} EPA issues \S~111(d) emissions guidelines for existing power plants based on control technology that each plant can install on site, so that every regulated plant can comply by operating more cleanly rather than by shifting generation elsewhere. That is the inside-the-fenceline, technology-based approach EPA had always used under \S~111(d), so the rule claims no ``unheralded'' power and the major questions doctrine likely does not apply. Challengers would instead have to argue under ordinary statutory standards---for example, that the technology is not ``adequately demonstrated'' or is unreasonably costly. \textit{West Virginia} itself expressly declined to decide whether \S~111(d) is limited to such measures.

\item \textbf{Fence-Sitter (Close Call): Mandatory Climate-Risk Disclosure by Public Companies.} Invoking its general authority to require disclosures ``necessary or appropriate in the public interest or for the protection of investors,'' the Securities and Exchange Commission requires public companies to report their greenhouse-gas emissions and climate-related financial risks. On one hand, disclosure rules sit within the agency's core expertise and long-standing practice, and the rule regulates what companies tell investors rather than how they produce energy. On the other hand, it imposes large compliance costs on a politically charged subject that Congress has not addressed, and it can be characterized as climate policy pursued through securities law. Whether a court treats the rule as an ordinary exercise of disclosure authority (as Justice Barrett's contextual reading of the doctrine in \textit{Biden v. Nebraska} might) or as a major question demanding clear authorization (as Justice Gorsuch's clear-statement approach might) could decide the case.

\end{enumerate}

\section{Critique}
\textbf{Substantive Canon vs. Contextual Textualism:} The Gorsuch--Barrett--Kagan debate across \textit{West Virginia}, \textit{Nebraska}, and \textit{Learning Resources} is a premier teaching vehicle for statutory interpretation: if a court demands ``clear congressional authorization'' even when a broad term (``system,'' ``modify,'' ``regulate'') textually covers the action, is the Court faithfully reading Congress's text or enforcing an unenumerated, judicially managed constitutional avoidance rule?

\textbf{Malleability of ``Majorness'' and Partisan Neutrality:} Critics initially charged that the major questions doctrine was a one-way deregulatory ratchet aimed at Democratic administrations (\textit{West Virginia}, \textit{NFIB}, \textit{Nebraska}). \textit{Learning Resources v. Trump} (2026) directly tests that critique---showing Roberts, Gorsuch, and Barrett applying the doctrine to strike down a Republican President's signature tariff policy, while Kavanaugh, Thomas, and Alito dissented.

\textbf{The ``Failed Legislation'' Trap:} Justice Kagan's critique of relying on failed subsequent bills highlights a classic textualist tension: congressional inaction can reflect filibuster gridlock or belief that the agency \textit{already} possesses statutory authority, not repudiation of an earlier Congress's enacted text.

\section{Key Quotations}
\begin{quote}
``Thus, in certain extraordinary cases, both separation of powers principles and a practical understanding of legislative intent make us `reluctant to read into ambiguous statutory text' the delegation claimed to be lurking there. To convince us otherwise, something more than a merely plausible textual basis for the agency action is necessary. The agency instead must point to `clear congressional authorization' for the power it claims.'' --- Chief Justice Roberts, \textit{West Virginia v. EPA}, 597 U.S. 697 (2022) (\href{https://supreme.justia.com/cases/federal/us/597/20-1530/}{Justia}).
\end{quote}
\begin{quote}
``But of course almost anything could constitute such a `system'; shorn of all context, the word is an empty vessel. Such a vague statutory grant is not close to the sort of clear authorization required by our precedents.'' --- Chief Justice Roberts (\href{https://supreme.justia.com/cases/federal/us/597/20-1530/}{Justia}).
\end{quote}
\begin{quote}
``Capping carbon dioxide emissions at a level that will force a nationwide transition away from the use of coal to generate electricity may be a sensible `solution to the crisis of the day.' But it is not plausible that Congress gave EPA the authority to adopt on its own such a regulatory scheme in Section 111(d). A decision of such magnitude and consequence rests with Congress itself, or an agency acting pursuant to a clear delegation from that representative body.'' --- Chief Justice Roberts (\href{https://supreme.justia.com/cases/federal/us/597/20-1530/}{Justia}).
\end{quote}
\begin{quote}
``When Congress seems slow to solve problems, it may be only natural that those in the Executive Branch might seek to take matters into their own hands. But the Constitution does not authorize agencies to use pen-and-phone regulations as substitutes for laws passed by the people's representatives.'' --- Justice Gorsuch, concurring (\href{https://supreme.justia.com/cases/federal/us/597/20-1530/}{Justia}).
\end{quote}
\begin{quote}
``The current Court is textualist only when being so suits it. When that method would frustrate broader goals, special canons like the `major questions doctrine' magically appear as get-out-of-text-free cards.'' --- Justice Kagan, dissenting (\href{https://supreme.justia.com/cases/federal/us/597/20-1530/}{Justia}).
\end{quote}

\end{document}
